\section{问题重述}
\subsection{研究背景}
小区微网通过光伏发电、储能电池和外部购电共同满足负荷。电池能够将低价时段的电能转移到高价时段，也能够吸收光伏富余电量；但预测偏差会使计划与实际需求不匹配，导致额外购电或弃电。因而，调度问题不仅是寻找一条低成本的购电曲线，还要使这条曲线在信息逐步到达的条件下可执行。
\subsection{四个问题的递进关系}
问题一依据典型日电价、负荷与光伏预测，建立当天购电及储能调度模型，要求日初、日末电量相同，并给出指定时段结果。问题二考虑逐日变化的负荷和光伏，要求每天零点作出购电计划，实际短缺按当期五倍电价紧急购电，常规购电按计划结算。
问题三允许利用零点及6、12、18时发布的光伏预报制定和调整计划，调整还会产生相对原计划的偏差费用，需要同时评价新增信息和调整时机的作用。
问题四进一步考虑波动电价，分别重算无日内调整与有日内调整两种策略。
本文在共同的储能约束下逐步改变信息集与结算目标，使四问结果具有清晰的比较基础。
